\documentclass[11pt]{article}
\usepackage{amsmath,amssymb,amsthm}
\usepackage[margin=1in]{geometry}
\usepackage{hyperref}

\newtheorem{theorem}{Theorem}
\newtheorem{corollary}[theorem]{Corollary}
\newtheorem{lemma}[theorem]{Lemma}
\theoremstyle{definition}
\newtheorem{definition}{Definition}
\theoremstyle{remark}
\newtheorem{remark}{Remark}

\title{A Self-Describing Cyclic Number:\\
The Unique Cyclotomic Decomposition of 142857\\
into Hamming-Family Scaffold Parameters}
\author{Vitex Probasco\\
\small University of Massachusetts Amherst (non-academic affiliate)\\
\small \texttt{vitexprobasco@umass.edu}}
\date{June 2026}

\begin{document}
\maketitle

\begin{abstract}
The cyclic number $142857 = (10^6-1)/7$ has prime-power factorization
$3^3 \cdot 11 \cdot 13 \cdot 37$. We show this equals
$r^3(n+k)(2n-1)(2^k+nr)$ where $r=3$, $n=7$, $k=4$ are the parameters
of the $[7,4,3]$ binary Hamming code and $b=10=n+r$ is the base.
We prove this cyclotomic--Hamming alignment is unique: among all binary
Hamming codes $[2^r-1,\, 2^r-1-r,\, r]$ with the canonical base
$b=n+r$, only $r=3$ produces a cyclic number whose cyclotomic
factor-blocks decompose into code parameters. The key uniqueness
condition, $2^{r-1}=r+1$, has the single positive integer solution
$r=3$. Two independent conditions from other cyclotomic factors
confirm this. A computational null test through $p=47$ verifies that
no other full-reptend prime in base~10 exhibits comparable
factorization alignment.
\end{abstract}

\section{Introduction}

The number $142857$ is the repetend of $1/7$ in base~10 and is
well-known as a cyclic number: multiplication by $1$ through~$6$
permutes its digits cyclically. Its prime factorization is
\[
142857 = 3^3 \cdot 11 \cdot 13 \cdot 37.
\]
We observe that each prime-power block admits a low-complexity
expression in the parameters of the $[7,4,3]$ Hamming code:
\begin{align}
27 &= 3^3 = r^3, \label{eq:27}\\
11 &= 7+4 = n+k, \label{eq:11}\\
13 &= 2\cdot 7-1 = 2n-1, \label{eq:13}\\
37 &= 2^4+7\cdot 3 = 2^k+nr. \label{eq:37}
\end{align}
Here $r=3$ (check/parity count), $n=2^r-1=7$ (code length),
$k=n-r=4$ (data width), and the canonical base $b=n+r=10$.
We use $\{n,k,r\}$ for the Hamming-family scaffold; in the
length-7 case this matches the familiar $[7,4,3]$ notation
because the minimum distance is also~$3$. Thus
\begin{equation}\label{eq:main}
142857 = r^3(n+k)(2n-1)(2^k+nr).
\end{equation}

This factorization is not a numerical coincidence imposed
post-hoc. It arises from the cyclotomic decomposition of
$b^{n-1}-1$ evaluated at $b=n+r$, and we prove the alignment
is unique to $r=3$.

\section{Cyclotomic Decomposition}

Since $n-1=6$, we have
\[
b^6-1 = \Phi_1(b)\,\Phi_2(b)\,\Phi_3(b)\,\Phi_6(b)
\]
where $\Phi_d$ denotes the $d$-th cyclotomic polynomial. At $b=10$:
\begin{align}
\Phi_1(10) &= 10-1 = 9 = 3^2, \\
\Phi_2(10) &= 10+1 = 11, \\
\Phi_3(10) &= 10^2+10+1 = 111 = 3 \cdot 37, \\
\Phi_6(10) &= 10^2-10+1 = 91 = 7 \cdot 13.
\end{align}
Therefore $10^6-1 = 9 \cdot 11 \cdot 111 \cdot 91 = 999999$,
and dividing by $n=7$:
\[
\frac{10^6-1}{7} = 3^3 \cdot 11 \cdot 13 \cdot 37 = 142857.
\]

Matching cyclotomic factors to Hamming parameters:
\begin{center}
\begin{tabular}{lll}
\hline
Cyclotomic factor & Value & Hamming expression \\
\hline
$\Phi_1(10) = 9$ & $3^2$ & $r^2$ \\
$\Phi_2(10) = 11$ & $11$ & $n+k$ \\
$\Phi_3(10) = 111$ & $3 \cdot 37$ & $r(2^k+nr)$ \\
$\Phi_6(10) = 91$ & $7 \cdot 13$ & $n(2n-1)$ \\
\hline
\end{tabular}
\end{center}
After division by $n=7$ (absorbing one factor from $\Phi_6$) and
collecting the factors of~$3$ ($r^2$ from $\Phi_1$ and $r$ from
$\Phi_3$), we obtain identity~\eqref{eq:main}.

\section{Uniqueness}

\begin{definition}
For $r \geq 2$, define the \emph{Hamming-family parameters}
$n = 2^r-1$, $k = n-r$, and the \emph{canonical base} $b = n+r$.
\end{definition}

We require $n$ to be prime (so that $\gcd(b,n)=1$ and $1/n$ has a
well-defined repetend in base~$b$). Define
$Q = (b^{n-1}-1)/n$, the Fermat quotient. When $\text{ord}_n(b)=n-1$,
$Q$ is the full-period cyclic repetend; otherwise it is the Fermat
quotient over the full exponent window.

\begin{theorem}\label{thm:main}
Among Hamming-family parameters with $n=2^r-1$ prime and $r \geq 2$,
the prime-power factor-blocks of $Q$ decompose into expressions of
$\{n,k,r\}$ only at $r=3$.
\end{theorem}

\begin{proof}
We establish three conditions on the cyclotomic factors. The first
alone suffices for uniqueness; the second and third provide
independent structural confirmation.

\medskip
\noindent\textbf{Condition 1 (uniqueness spine).}
We require $\Phi_2(b) = b+1 = n+k$. Since $b=n+r$ and $k=n-r$:
\[
n+r+1 = n + (n-r) = 2n-r,
\]
which simplifies to $n = 2r+1$, equivalently
\begin{equation}\label{eq:cond1}
2^{r-1} = r+1.
\end{equation}
At $r=3$: $2^2 = 4 = 4$. \checkmark

For $r \geq 4$: $f(r) = 2^{r-1}-(r+1)$ satisfies $f(4)=3>0$, and
$f(r+1)-f(r) = 2^{r-1}-1 > 0$ for $r \geq 2$, so $f$ is strictly
increasing for $r \geq 3$. Thus $r=3$ is the unique solution.

\medskip
\noindent\textbf{Condition 2 (independent confirmation).}
We require $n \mid \Phi_6(b)$ with $\Phi_6(b)/n = 2n-1$.
Since $\Phi_6(b) = b^2-b+1 = (n+r)^2-(n+r)+1$, expanding and
setting equal to $n(2n-1)$ yields
\begin{equation}\label{eq:cond2}
n(n-2r) = r^2-r+1.
\end{equation}
At $r=1$: $1 \cdot (1-2) = -1 \neq 1$. At $r=2$: $3 \cdot (3-4) = -3 \neq 3$.
At $r=3$: $7 \cdot (7-6) = 7 = 7$. \checkmark

For $r \geq 4$: the left side grows as $\Theta(4^r)$ while the right
grows as $\Theta(r^2)$, so no further solutions exist.

\medskip
\noindent\textbf{Condition 3 (for prime $r$).}
We require $r \mid \Phi_3(b)$ where $\Phi_3(b) = b^2+b+1$.
Since $b \equiv n \pmod{r}$ and $n = 2^r-1$:
\[
\Phi_3(b) \equiv n^2+n+1 = 2^{2r}-2^r+1 \pmod{r}.
\]
By Fermat's little theorem (for prime~$r$): $2^r \equiv 2 \pmod{r}$,
so $2^{2r} \equiv 4 \pmod{r}$, giving
\[
\Phi_3(b) \equiv 4-2+1 = 3 \pmod{r}.
\]
Thus $r \mid \Phi_3(b)$ requires $r \mid 3$, and for prime $r \geq 2$
this gives $r=3$.

\begin{remark}
Condition~3 uses Fermat's little theorem and therefore applies only
to prime~$r$. In the theorem's domain, $n=2^r-1$ prime implies $r$
prime (contrapositive: $r$ composite implies $2^r-1$ composite), so
the Fermat step is licensed. Composite solutions to the divisibility
condition alone exist (e.g., $r=57$), but fall outside the hypothesis.
In any case, Condition~1 already establishes uniqueness unconditionally.
\end{remark}

The unique simultaneous solution is $r=3$, giving the $[7,4,3]$
Hamming code in base $b=10=7+3$.
\end{proof}

\begin{corollary}
Base~$10 = n+r$ is the canonical-base candidate for the $\{7,4,3\}$
scaffold. More generally, each Hamming-family code has canonical-base
candidate $b = 2^r-1+r$, but only at $r=3$ does the resulting
Fermat quotient exhibit the self-describing factorization.
\end{corollary}

\section{Supplementary Observations}

\subsection{Half-period split (Midy-type)}
The repetend splits as $142\,|\,857$ with
\[
142+857 = 999 = 10^3-1
\]
(a consequence of Midy's theorem for even-period primes) and
\[
857-142 = 715 = 5 \cdot 11 \cdot 13.
\]
The difference factors as $5 \cdot (n+k) \cdot (2n-1)$, coupling the
base-10 cofactor ($5$, since $10=2\times5$) with two scaffold primes.

The split has a compact derivation. Let $Q = (10^3+1)/7 = 143 = 11
\times 13$. Then $142 = Q-1$ and $857 = 1000-Q$, so:
\[
857-142 = 1001 - 2Q = 5Q = 5 \times 11 \times 13.
\]
The same residual~$5$ appears in both the split difference and the
half-cycle correction below.

\subsection{Half-cycle correction}
\[
\frac{1}{7} - \frac{142}{999} = \frac{5}{6993} = \frac{5}{7 \times 999}.
\]
The correction from half-cycle repetition to the exact value is the
base cofactor divided by the code length times the nines complement.

\subsection{Digit invariants}
The digit set $\{1,2,4,5,7,8\}$, invariant under cyclic rotation, has
\[
\text{sum} = 27 = 3^3 = r^3, \qquad
\text{product} = 2240 = 2^6 \cdot 5 \cdot 7.
\]

\section{Computational Null Test}

We verified that no other full-reptend prime $p \leq 47$ in base~10
produces a cyclic number whose complete factorization decomposes into
expressions of its own structural parameters. Additionally, $p=31$
(the next Mersenne-prime Hamming length, $31=2^5-1$) is not
full-reptend in base~10: its period is~15, not~30. See the
accompanying verification code.

\begin{thebibliography}{9}
\bibitem{hamming} R.~W.~Hamming, ``Error Detecting and Error
  Correcting Codes,'' \emph{Bell System Technical Journal}, vol.~29,
  no.~2, pp.~147--160, 1950.

\bibitem{midy} E.~Midy, ``De Quelques Propri\'et\'es des Nombres
  et des Fractions D\'ecimales P\'eriodiques,'' Nantes, 1836.

\bibitem{nz} I.~Niven, H.~S.~Zuckerman, and H.~L.~Montgomery,
  \emph{An Introduction to the Theory of Numbers}, 5th ed., Wiley, 1991.
\end{thebibliography}

\end{document}
